\corrigesChapitre{2}{Exponentielle et logarithme}
\begin{corrige}{1}{Identités immédiates}
\begin{questions}
\item $0$, $3$ et $5$, par les identités réciproques.
\item $\ln(8/2)=\ln4=2\ln2$.
\item $9\sqrt3=3^{5/2}$ donc $\ln(9\sqrt3)=\tfrac52\ln3$.
\end{questions}
\end{corrige}
\begin{corrige}{2}{Domaine avant calcul}
\begin{questions}
\item $2x-6>0$, soit $x>3$.
\item $x^2>1$, soit $]-\infty,-1[\cup]1,+\infty[$.
\item Non : la première exige $x>0$ ; la seconde exige seulement $x\ne0$.
\end{questions}
\end{corrige}
\begin{corrige}{3}{Équations simples}
\begin{questions}
\item $2x=\ln7$, donc $x=\tfrac12\ln7$.
\item Sur $x>-2$, $x+2=1$, donc $x=-1$.
\item Sur $x>0$, la croissance de $\ln$ donne $0<x\leq\mathrm e^2$.
\end{questions}
\end{corrige}
\begin{corrige}{4}{Dériver}
\begin{questions}
\item $x>-1/3$ et $f\,\!^{\prime}(x)=3/(3x+1)$.
\item $g\,\!^{\prime}(x)=-2\mathrm e^{-2x}$.
\item $x^2+1>0$ sur $\R$, donc $h\,\!^{\prime}(x)=2x/(x^2+1)$.
\end{questions}
\end{corrige}
\begin{corrige}{5}{Un seuil de croissance}
\begin{questions}
\item $n>\ln(1{,}6)/\ln(1{,}03)\simeq15{,}901$.
\item Le premier rang est $16$ ; $u_{15}\simeq389{,}49$ et $u_{16}\simeq401{,}18$.
\item L’inégalité sur $n$ deviendrait large ; ici le quotient n’étant pas entier, le rang resterait $16$.
\end{questions}
\end{corrige}
\begin{corrige}{6}{Une décroissance}
\begin{questions}
\item $T=\ln2/0{,}2\simeq3{,}466$ heures.
\item $\mathrm e^{-0{,}2t}<1/8$ équivaut à $t>\ln8/0{,}2\simeq10{,}397$ heures.
\item $C(t)\to0$ ; le modèle décrit une approche de zéro, pas une annulation exacte en temps fini.
\end{questions}
\end{corrige}
\begin{corrige}{7}{Somme de logarithmes}
\begin{questions}
\item Le domaine est $x>2$. On obtient $x(x-2)=3$, soit $(x-3)(x+1)=0$, donc $x=3$.
\item $-1$ annule le polynôme mais n’appartient pas au domaine des logarithmes.
\item Sur $x>2$, $x^2-2x-3\leq0$ donne $2<x\leq3$.
\end{questions}
\end{corrige}
\begin{corrige}{8}{Changer de base}
\begin{questions}
\item $3$ et $5$, car $1\,000=10^3$ et $32=2^5$.
\item $\log_2 x=\ln x/\ln2$ pour $x>0$.
\item $\log_{10}(2^{20})=20\log_{10}2\simeq6{,}0206$ ; le nombre est compris entre $10^6$ et $10^7$ et possède donc $7$ chiffres.
\end{questions}
\end{corrige}
\begin{corrige}{9}{Une propriété à prouver}
\begin{questions}
\item $a(1/a)=1$, donc $\ln a+\ln(1/a)=\ln1=0$.
\item $\ln(a/b)=\ln(a\times1/b)=\ln a-\ln b$ pour $a,b>0$.
\item Avec $a=3,b=1$, le premier membre vaut $\ln2$ et le second $\ln3$ : ils sont différents.
\end{questions}
\end{corrige}
\begin{corrige}{10}{Tables de calcul}
\begin{questions}
\item $\ln6\simeq1{,}791759$ ; $\ln(3/2)\simeq0{,}405465$.
\item L’inégalité triangulaire donne dans les deux cas une erreur au plus $10^{-6}$.
\item Les produits et quotients deviennent des sommes et différences ; il faut cependant une lecture inverse de table et maîtriser les arrondis.
\end{questions}
\end{corrige}
\begin{corrige}{11}{Comparer deux modèles}
\begin{questions}
\item $A\,\!^{\prime}(0)=20$, $B\,\!^{\prime}(0)=10$. Le modèle affine commence plus vite.
\item Diviser par $100$ donne $1+0{,}2t=\mathrm e^{0{,}1t}$ ; les deux membres sont positifs et l’application de $\ln$ donne l’équivalence.
\item Non : $t$ reste des deux côtés. On utilise une étude de fonction puis une méthode numérique ; $t=0$ est déjà une solution.
\end{questions}
\end{corrige}
\begin{corrige}{12}{Le piège des inégalités}
\begin{questions}
\item $n\geq\ln(1/40)/\ln(0{,}9)\simeq35{,}012$ ; le premier rang est $36$.
\item $\ln(0{,}9)<0$ : diviser par ce nombre inverse le sens.
\item Initialiser $n=0,v=40$ et répéter tant que $v>1$ : remplacer $v$ par $0{,}9v$ et augmenter $n$ de $1$.
\end{questions}
\end{corrige}
